\begin{frame}{Czym się zajmujemy?}
    % \textcolor{cfblue}{Dychotomia CSP (Bulatov, Zhuk; 2017)} \\
    % Dla grafu \(H\) problem CSP\((H)\) albo jest w P, albo NP-zupełny.

    % \vspace{1em}

    \textcolor{cfblue}{Hipoteza (Bang-Jensen i Hell, 1990)} \\
    Niech \(H\) będzie spójnym grafem bez źródeł i ujść. Jeśli \(H\) ma homomorfizm do cyklu, to~problem CSP\((H)\) jest w klasie P, wpp CSP\((H)\) jest NP-zupełny.

    \vspace{1em}
\begin{table}[]
    \centering

\begin{tabular}{|c|c|c|}
\hline
Kroki & Dla rdzenia & Uogólnienie  \\
\hline
Relacja centralna & \ding{51} & \textcolor{cfblue}{\ding{51}} \\
\hline
DorD dla CE & \ding{51} & \textcolor{cfblue}{\ding{51}} \\
\hline
DorD & \ding{51} &  \\
\hline
=or= & \ding{51} &  \\
\hline
NP-zupełność & \ding{51} &  \\
\hline
\end{tabular}

\end{table}
\end{frame}